\section{Conclusions}
\label{rz:conclusiones}

La forme de Weil admet, sur chaque fenêtre fixée, une réduction exacte à un problème fini qui conserve son complément infini. L'espace des détails de moyenne nulle peut être rendu coercif par raffinement ; son élimination produit un complément de Schur avec un critère équivalent de positivité et le même indice négatif que la forme complète. Dans la fenêtre $(-1/18,1/18)$, le contrôle conjoint des moyennes, des détails et des couplages démontre
\[
\mathscr W(f,f)\geq\frac12\|f\|_2^2,
\qquad f\in V_{1/9}.
\]
L’affirmation comprend la fermeture des fonctions tests lisses
à support compact dans la norme conjointe des lecteurs.
C’est un résultat sur un domaine fonctionnel infini, stable
par translation et raffinement intérieur.

\subsection*{Signe de la forme et conservation du complément infini}

L’identité
\[
\mathscr W_R(Ua+d,Ua+d)
=a^*\bigl(A-B^*D^{-1}B\bigr)a
+\mathfrak d[d+D^{-1}Ba,d+D^{-1}Ba]
\]
explique ce que conserve la réduction. Le second terme est
positif et contient tous les détails ; le premier conserve
leur interaction avec les moyennes au moyen de $B^*D^{-1}B$.
Le théorème~\ref{schur-add-schur-exacto} prouve que la
positivité sur $V_R$ équivaut à celle de cette matrice.
Sa taille finie ne remplace pas le calcul de ses entrées :
la certification résiduelle de
\eqref{schur-add-eq-S7} fournit des bornes avec l'erreur du
complément retenu.

Dans la fenêtre de longueur $1/9$, les bornes $D>I$, $A>97/100$ et $\|B\|^2<1/5$ produisent la marge coercive annoncée. Ce cas local ne contient aucune translation première active puisque $1/9<\log2$. L'identité de transport de la forme complète permet toutefois de porter la même borne sur chaque sous-espace $C_a^n\tau_t C_c^\infty((-1/18,1/18))$, avec $a>b>1/2$. Ses éléments peuvent avoir des queues exponentielles, et leur contribution première est conservée avec la somme convergente complète. L'estimation s'applique à chaque collection entière, et non à une seule fonction test sélectionnée.

Les autres domaines positifs maintiennent leur contenu propre : l'assemblage spécifié de neuf intervalles contrôle leurs produits croisés et tous leurs raffinements, tandis qu'une fenêtre arbitraire admet un seuil de coercivité pour les détails cellulaires. La distinction entre ces domaines permet d'identifier aussi bien la marge démontrée que le couplage qui décide du signe d'une composition plus large.

\subsection*{L'arithmétique détermine les modes et leurs lecteurs}

L'antécédent arithmétique est construit avant la lecture
spectrale. Les deux feuillets d'APP fournissent les opérations
locales et leurs quotients ; le produit natif propage ces données
dans des mots de forme normale unique. Les ouvertures du produit
sélectionnent les irréductibles, et l'évaluation les identifie aux
nombres premiers positifs. La réalisation opératorielle du sélecteur
compte les décompositions non triviales. Sa positivité est
une norme au carré avec cette signification arithmétique précise.

La factorisation organise ensuite les occupations des modes et
leurs horloges. Le produit d’Euler et les moments
premier--puissance proviennent des traces de cette représentation.
La limite de la trace de chaleur cyclique et sa transformation de
Mellin produisent la fonction complétée, avec ses facteurs gamma
et polaires et son équation fonctionnelle. Les parties finies, les
évaluations régularisées et les bornes de reste conservent
les normalisations nécessaires pour composer ces lectures.

L'unité du parcours réside dans les états enrichis d'APP--TRIT--TPK et dans la structure conjointe du continu : route, orientation, feuillets, mémoire et incidences sont transmis à travers des restrictions et des raffinements compatibles. L'évaluation réelle obtient un résultat de cette construction ; l'histoire et la frontière conservent l'information dont leurs réalisations ultérieures ont besoin.

Les lois de composition précisent cette continuité. Le produit de
deux moments se réalise sur des paires de modes et leurs coefficients se
regroupent par convolution des diviseurs ; la composition diagonale maintient
un seul mode. La normalisation spectrale transporte tout le développement
de Laurent, y compris le terme polaire, par une transformation
triangulaire affine inversible. L'équation fonctionnelle relie les
dérivées négatives paires aux valeurs impaires positives et
établit en particulier $Z(3)=4\pi^2\log A_2$. Enfin, les coupures
par irréductibles apportent des encadrements explicites complémentaires aux
coupures par profondeur. Chaque identité conserve l'opération qui la
produit et le domaine dans lequel elle est démontrée.

\subsection*{Frontière, vacances et double cercle}

La reconstruction de l'intérieur à partir de la frontière relevée dans la coordinatisation multiaffine fournit une instance explicite de conservation. La différence somme--produit conserve la donnée résiduelle et le quotient ; les lectures $72,27,77,22$ et la complétion $729+271$ maintiennent leurs opérations et leurs domaines. Le centre $(5,5)$ et ses vacances conjuguées conservent la somme et le résidu multiplicatif tandis que le quotient change. La signature régionale $555555$ requiert trois feuillets de continuation, et le calendrier de capacité produit les longueurs $21,22$ et les séparations $6,7$ avec un bilan exact.

Le double cercle porte cette conservation au transport des lecteurs. La compression de neuf phases produit le complément $(8z,-z,\ldots,-z)/27$ et le facteur $8/81$. Le rapport modal $q=5/9$ détermine les moments de mémoire $q^{|n|}$ ; leur factorisation triangulaire prouve la positivité en toute dimension et le déterminant $(56/81)^N$. Le registre des pertes conserve le terminal et les produits croisés. Le raffinement peut conduire à une limite faible atomique d'états, dont la représentation se distingue du décalage régulier dans lequel les vecteurs initiaux ont été construits.

Pour les lecteurs arithmétiques, la queue gamma reconstruit la
fonction test et les lignes des deux colonnes. Le connecteur résultant
conserve exactement
$J_+^*J_+-J_-^*J_-$, avec les termes premier, gamma et polaire.
Son action permet de formuler le problème énergétique sur un
opérateur explicite. L'invariance sous $C_a$ préserve en outre
tous les moments croisés de la forme, et pas seulement ses
évaluations diagonales.

\subsection*{Condition globale et relation avec l’hypothèse de Riemann}

Le critère global de Weil exige la positivité sur tout le domaine des fonctions tests. Dans la formulation des deux colonnes, la condition exacte est le caractère contractant du transfert atteignable sur l'adhérence de l'image positive, avec compatibilité lors de l'élargissement des fenêtres. Dans la formulation circulaire, elle s'exprime par l'identité
\[
m_n^{\rm ar}(f,g)
=\langle U_{\rm mem}^{\,n}\mathcal E f,
          \mathcal E g\rangle_{\mathcal K_{\rm lim}},
\]
pour tous les moments mixtes et une collection couvrant le domaine au moyen des limites déclarées. L'identité d'énergie et l'entrelacement doivent être établis conjointement ; la positivité du Gram régulier de mémoire ne les remplace pas.

Les preuves de cet article établissent la reconstruction, l'invariance et la coercivité sur les domaines indiqués. L'identification globale précédente reste une hypothèse des résultats spectraux qui l'emploient et ne se déduit pas de la positivité de chaque collection transportée séparément. L'article ne conclut donc pas à une résolution de l'hypothèse de Riemann. Sa contribution fonctionnelle est précise : il transforme la conservation de la structure arithmétique et de la mémoire en opérateurs récupérables, en réductions exactes et en bornes coercives, en maintenant visibles les produits mixtes et les domaines qui gouvernent tout prolongement global.
